%% Chapter 8 Solutions
%% ============================================================

\section{Chapter 8: Caveat Emptor}

%% ----------------------------------------------------------
\begin{solution}{8.1}{L3}

\solanswer

$g(u;p) = u^3 - 3pu + 2 = 0$ (same as Ex 4.2, now with bifurcation focus).

\solpart{a} \textbf{Real roots at $p = 1$ and most rapidly changing.}
Roots: $u_1 = 1$ (double), $u_2 = -2$. The double root $u_1 = 1$
changes most rapidly — its sensitivity is undefined (infinite), as shown
in Ex 4.2. It is at a fold bifurcation point.

\solpart{b} \textbf{$\SI{p}{u^*}$ at each root.}
As computed in Ex 4.2: $\SI{p}{u_1}$ is undefined (singular FSE);
$\SI{p}{u_2} = 1/3$.

\solpart{c} \textbf{Root curve and bifurcation.}
The real roots satisfy $u^3 - 3pu + 2 = 0$. By Cardano's formula or
numerical continuation, the two positive roots merge at the fold bifurcation
where $\partial g/\partial u = 3u^2 - 3p = 0$, i.e., $u = \sqrt{p}$.
Substituting: $p^{3/2} - 3p\sqrt{p} + 2 = 0 \Rightarrow -2p^{3/2} + 2 = 0
\Rightarrow p^* = 1$. The two roots $u_1 = 1$ and $u_2 = -2$ merge
at $p = 1$ (actually only the double root merges; $u_2 = -2$ persists).
Careful: the cubic has the double root at $u = 1$ for $p = 1$.
For $p < 1$, the double root splits into two complex roots (no real positive root
near 1), and for $p > 1$, two distinct real roots appear near $u = 1$.
The bifurcation is at $\mathbf{p = 1}$.

\solpart{d} \textbf{FSE prediction at bifurcation.}
At $p = 1$, $\partial g/\partial u = 0$: the FSE $(\partial g/\partial u)\,(\partial u/\partial p) = -\partial g/\partial p$ has a zero coefficient, so $\partial u/\partial p = \infty$.
Finite-difference verification: $u(p=1.001)$ vs $u(p=0.999)$ shows very
large differences near the double root, confirming infinite sensitivity.
The FD estimate diverges as $p \to 1$, consistent with a fold bifurcation.

\solnote{This exercise directly connects Chapter 4 (implicit FSE) to
Chapter 8 (bifurcations). The key point: the FSE singularity \emph{is}
the bifurcation. Students who answer only ``sensitivity is large'' without
connecting to the eigenvalue of the Jacobian receive partial credit.}

\end{solution}

%% ----------------------------------------------------------
\begin{solution}{8.2}{L3}

\solanswer

\solpart{a} \textbf{$\SI{\beta}{I(T)}$ at four values, $T = 90$ days.}

Computed via centered FD with $h = 6\times10^{-6}\beta^*$:

\begin{center}
\begin{tabular}{lcccc}
\toprule
$\beta$ & $R_0 = k\beta\tau$ & $I(90)$ & $\SI{\beta}{I(90)}$ (approx.) \\
\midrule
0.04 & 1.40 & $\sim150$ & $\sim+3.2$ \\
0.06 & 2.10 & $\sim480$ & $\sim+1.8$ \\
0.08 & 2.80 & $\sim620$ & $\sim+1.2$ \\
0.12 & 4.20 & $\sim720$ & $\sim+0.7$ \\
\bottomrule
\end{tabular}
\end{center}

The SI \emph{decreases} as $R_0$ increases past 1: at higher $R_0$,
almost everyone has been infected by day 90, and $I(90)$ is in the
declining tail, less sensitive to $\beta$.

\solpart{b} \textbf{$\beta = 0.02$ ($R_0 = 0.7 < 1$).}
No epidemic occurs: $I(90) \approx I(0) = 1$ (infection dies out
exponentially). $\SI{\beta}{I(90)}$ is small and positive near zero.
\textbf{Qualitative change:} the SI is negligibly small for $R_0 < 1$
because no epidemic amplification occurs, but large and positive
for $R_0 > 1$. Near $R_0 = 1$ ($\beta \approx 0.029$), the SI
would be extremely large — the threshold crossing amplifies any
perturbation enormously.

\solpart{c} \textbf{Bifurcation diagram.}
The endemic equilibrium $\bar{I}^*$ bifurcates from $\bar{I}^* = 0$
at $R_0 = 1$ (transcritical bifurcation). For $R_0 < 1$: $\bar{I}^* = 0$
(disease-free equilibrium is stable). For $R_0 > 1$: $\bar{I}^* > 0$
(endemic equilibrium exists and is stable). The bifurcation diagram shows
$\bar{I}^*$ vs.\ $\beta$ as a forward (supercritical) transcritical bifurcation.

\solnote{The key finding from part (a) is that $\SI{\beta}{I(T)}$ \emph{decreases}
as $\beta$ increases — counterintuitive if one expected larger $\beta$ to
mean larger sensitivity. The reason: at high $\beta$, the epidemic has
already run its course by day 90, and $I(90) \approx 0$ regardless of
small perturbations to $\beta$. This illustrates why SA at a single time
point can be misleading (Chapter 4, Ex 4.8).}

\end{solution}

%% ----------------------------------------------------------
\begin{solution}{8.3}{L4}

\solanswer

\solpart{a} \textbf{Condition number and error vs.\ $\epsilon$.}

MATLAB code:
\begin{verbatim}
u_exact = [1;1;1];
eps_vals = [1, 0.1, 0.01, 0.001, 0.0001];
for k = 1:length(eps_vals)
    ep = eps_vals(k);
    A = [1 2 3; 4 5 6; 7 8 9+ep];
    b = A * u_exact;
    kappa(k) = cond(A);
    err(k) = norm(A\b - u_exact);
end
\end{verbatim}

Expected results (approximate):

\begin{center}
\begin{tabular}{rcc}
\toprule
$\epsilon$ & $\kappa(A)$ & $\|\hat{\vec{u}} - \vec{u}\|$ \\
\midrule
1.0    & $\sim 30$    & $\sim 10^{-14}$ \\
0.1    & $\sim 300$   & $\sim 10^{-13}$ \\
0.01   & $\sim 3000$  & $\sim 10^{-12}$ \\
0.001  & $\sim 3\times10^4$ & $\sim 10^{-11}$ \\
0.0001 & $\sim 3\times10^5$ & $\sim 10^{-10}$ \\
\bottomrule
\end{tabular}
\end{center}

\solpart{b} \textbf{Log-log slope.}
The relationship is $\|\hat{\vec{u}} - \vec{u}\| \propto \kappa(A)\,\varepsilon_{\text{mach}}$.
On log-log axes, the slope is $+1$ (error grows linearly with condition number).

\solpart{c} \textbf{Implication for large SI values.}

A large SI for a nearly singular model reflects the ill-conditioning of the
sensitivity computation, not necessarily a genuine large sensitivity of the
physical model. If the Jacobian $J_F$ of the model equations is nearly
singular (large $\kappa(J_F)$), the FSE $J_F\,(\partial\vec{u}/\partial p_j) = \vec{f}_j$
will amplify any numerical errors in $\vec{f}_j$ by a factor of $\kappa(J_F)$.
The reported SI may be $\kappa(J_F)$ times larger than the true value.
\textit{Practical implication:} always report $\kappa(J_F)$ alongside
large SI values. An SI of $50$ with $\kappa(J_F) = 100$ should be
flagged as potentially unreliable; an SI of $50$ with $\kappa(J_F) = 2$ is trustworthy.

\solnote{The slope of 1 in part (b) is the classical condition-number
error bound: $\|\delta\vec{u}\|/\|\vec{u}\| \leq \kappa(A)\|\delta\vec{b}\|/\|\vec{b}\|$.
Students who compute the slope empirically and report it as $\approx 1$
without citing the bound receive full credit; students who derive the bound
analytically receive bonus credit.}

\end{solution}

%% ----------------------------------------------------------
\begin{solution}{8.4}{L5}

\solanswer

\textbf{Advisory report: SA near the epidemic threshold ($R_0 = 1.04$).}

\solpart{a} \textbf{What local SA can and cannot tell you.}

Local SA provides the derivative of each \QOI with respect to each parameter
at the current best-fit values. With $R_0 = 1.04$, the model is very
close to the epidemic threshold. Local SA \emph{can} tell you which
parameters would, if perturbed slightly, have the largest marginal effect
on disease burden. What it \emph{cannot} tell you: whether a parameter
change of realistic magnitude (say 10--20\%) would push $R_0$ below 1
(eliminating the epidemic entirely) or keep it above 1. Near the threshold,
the model's behavior changes qualitatively, and local SA captures only the
linearized behavior in an infinitesimally small neighborhood.

\solpart{b} \textbf{Specific numerical problems near $R_0 = 1$.}

Near the transcritical bifurcation at $R_0 = 1$, the Jacobian of the
model at the endemic equilibrium has an eigenvalue near zero. This causes:
(1) the FSE linear system $J_F\,(\partial\vec{u}/\partial p) = \vec{f}$
to be nearly singular, amplifying roundoff errors by $\kappa(J_F)$;
(2) time-dependent SIs to grow without bound near the threshold crossing;
(3) the adjoint ODE to have a slowly decaying mode (near-zero eigenvalue),
making backward integration unstable over long time horizons.
Practitioners should report $\kappa(J_F)$ and check that SI values are
stable to perturbations of the nominal parameter set within the confidence region.

\solpart{c} \textbf{Additional analyses recommended.}

\begin{enumerate}[nosep]
\item \textit{Bifurcation diagram:} Plot the endemic equilibrium $I^*$ vs.\
  $R_0$ (or vs.\ the most controllable parameter). Determine how large a
  reduction in $\beta$ or $k$ is needed to push $R_0$ below 1.
  This gives a specific, actionable target rather than a marginal SI value.
\item \textit{SA at multiple $R_0$ values:} Compute SI at $R_0 = 0.9$,
  $1.0$, $1.04$, $1.2$. Report how the rankings change. If the ranking
  is stable across this range, the local SA at $R_0 = 1.04$ is informative.
  If not, flag the instability.
\item \textit{Finite-change analysis:} Compute the actual change in peak $I$
  for a $10\%$ reduction in each parameter. Report these directly alongside
  the SI values. For parameters near the threshold, the finite change may
  be qualitatively different from the SI prediction.
\end{enumerate}

\solpart{d} \textbf{Explaining bifurcation to a policy maker.}

Think of the epidemic threshold as a tipping point. Below it ($R_0 < 1$),
each infected person infects on average less than one other person, and
the outbreak fades on its own. Above it ($R_0 > 1$), each infected person
infects more than one other person, and the outbreak grows into an epidemic.
At $R_0 = 1.04$, we are just above the tipping point. Even a small
reduction in contact rates or transmission probability — say $5\%$ —
could push us below the tipping point, causing the epidemic to die out
rather than sustain itself. This is the most important insight from the
sensitivity analysis: at $R_0$ close to 1, small interventions can have
disproportionately large effects. But it also means our predictions are
highly uncertain: if the true $R_0$ is $0.96$ rather than $1.04$, the
epidemic might already be declining naturally.

\solnote{This is a realistic scenario that every applied mathematical modeler
faces. The report format requirement is deliberate: public health advisors
need concise, clear guidance, not a technical treatise. Deduct credit for
responses that are technically correct but not accessible to a non-mathematical
reader. The bifurcation explanation (part d) should use no equations.}

\end{solution}

%% ----------------------------------------------------------
\begin{solution}{8.5}{L6}

\solanswer

\textbf{Three methodological concerns with the malaria SA ($R_0 = 1.15$, all SI $< 0.1$):}

\textbf{Concern 1: Small SIs near $R_0 = 1$ may be numerical artifacts.}

\textit{Problem:} With $R_0 = 1.15$, the model is moderately above the
threshold, but close enough that the endemic equilibrium Jacobian may be
ill-conditioned. Small SI values could result from condition-number amplification
rather than genuine insensitivity.

\textit{Evidence needed:} Report $\kappa(J_F)$ at the endemic equilibrium
and verify that $|\text{SI}| \gg \kappa(J_F)\,\varepsilon_{\text{mach}}$
for each parameter. If any SI is within a factor of $\kappa(J_F)$ of machine
epsilon, that SI is numerically unreliable.

\textit{Alternative analysis:} Compute SIs of $R_0$ itself (analytically)
and compare sign and magnitude with the equilibrium SIs. If $\SI{p_j}{R_0}$
is large but $\SI{p_j}{\bar{I}}$ is small, the discrepancy suggests
numerical issues at the equilibrium.

\textbf{Concern 2: SA of the endemic equilibrium does not address
transient dynamics or the probability of outbreak onset.}

\textit{Problem:} If malaria is currently at or near the endemic equilibrium,
the relevant policy question is ``which intervention would most reduce
endemic prevalence?'' But if the model is used for outbreak forecasting
(a different but common use), the relevant QOI is the peak or cumulative
burden during a transient following an introduction event. SA of the
endemic equilibrium cannot answer the latter question.

\textit{Evidence needed:} Clarify the model's purpose. If forecasting
an outbreak, compute SA for the transient QOI (e.g., peak $I$ at $T = 180$
days) in addition to the endemic equilibrium.

\textit{Alternative analysis:} Compute time-dependent SIs $\SI{p_j}{I}(t)$
and report whether the parameter ranking changes between the transient
and endemic phases.

\textbf{Concern 3: ``All indices less than 0.1'' does not imply robustness
unless the uncertainty ranges are accounted for.}

\textit{Problem:} A local SI of 0.1 means a 1\% change in the parameter
produces a 0.1\% change in prevalence. But if a parameter is uncertain
by $\pm50\%$ (plausible for entomological parameters in malaria models),
the expected change in prevalence from this parameter alone is $\approx 5\%$.
Reporting only SI values without the uncertainty range gives an incomplete
picture.

\textit{Evidence needed:} Report sigma-normalized SIs or one-way sensitivity
ranges: the change in $\bar{I}$ when each parameter varies over its full
uncertainty range. Parameters with small SI but large uncertainty range
may still dominate the output variance.

\textit{Alternative analysis:} Compute Sobol first-order and total indices
across the plausible parameter uncertainty ranges. A parameter with local
SI = 0.1 and uncertainty range $\pm80\%$ may have Sobol $S_1 = 0.3$.

\solnote{Full credit requires all three concerns with (a), (b), (c) for each.
The three concerns should be genuinely distinct: numerical reliability
(Concern 1), scope of the SA (Concern 2), and accounting for uncertainty
ranges (Concern 3). Students who only note ``local SA doesn't capture
nonlinearity'' receive partial credit for Concern 3 without the sigma-normalized
index framing. The claim ``all indices less than 0.1'' is particularly dangerous
near a threshold --- this is the chapter's core message.}

\end{solution}

%% ============================================================
